% ---------------------------------------------
% Copyright Ben Paul Wise. All rights reserved.
% ---------------------------------------------
% process with TeXworks as pdfLaTeX+MakeIndex+BibTex


% This file is \input by main.tex as the STYLE MODULE, so the class line is
% commented out here and the \documentclass (report, for chapters) lives in
% main.tex; everything below applies unchanged.
%\documentclass[english, 10pt]{article}
% `english' is necessary for babel, not `letter' or `a4paper'
% default is 10pt unless specified otherwise.
% I reserve 12 pt for legible final documents


%\usepackage[a4paper]{geometry}
\usepackage[letterpaper]{geometry}

\usepackage{geometry}
\geometry{verbose,lmargin=2.0cm,rmargin=2.0cm}

\setlength{\parindent}{2em}
\setlength{\parskip}{0.75em}

\usepackage{fancyhdr}
\pagestyle{fancy}
\usepackage{color}
\definecolor{mypurple}{RGB}{128, 0, 128}
\definecolor{myTeal}{RGB}{0, 128, 128}
\definecolor{myDarkGreen}{RGB}{0, 92, 0}
\definecolor{myLightGreen}{RGB}{0, 192, 0}
\definecolor{myDarkRed}{RGB}{192, 0, 0}
\definecolor{myDarkBlue}{RGB}{0, 0, 192}
\definecolor{myBeige}{RGB}{255, 255, 192}

\definecolor{GroupWBlue1}{RGB}{75, 99, 175} 
\definecolor{GroupWGrey1}{RGB}{145, 147, 150} 

\definecolor{schrift}{RGB}{0,73,114}

\usepackage{amsmath} %% get AMS standards for theorems, e.g. the hollow QED box,  \qedsymbol
%% But I like the solid "tombstone", Halmos, whatever:
\newcommand{\qedsolid}{\rule{1ex}{1.62ex}}
\usepackage{amssymb} % \sum


\usepackage{empheq} % colored equation-boxes
\usepackage[most]{tcolorbox}
\tcbset{colback=white!90!yellow, % i.e. 90% and 10% yellow
        colframe=green!36!blue, % i.e. 64% green and 36% blue, `schrift'
        highlight math style= {enhanced, %<-- needed for the ’remember’ options
            colframe=red,colback=red!10!white,boxsep=0pt}
        }
		
		
% setup colored and boxed quotes	
% this would give a yellowish background inside a reddish frame:
% colback=yellow!10!white, colframe=red!50!black, 

%\usepackage[most]{tcolorbox}     
\newtcolorbox{myquote}{
colback=yellow!10!white,
colframe=blue!50!black,
grow to right by=-20mm,
grow to left by=-30mm,
%boxrule=2pt,
boxsep=0pt,
breakable} \makeatletter

\newcommand{\blockquote}[1]{  \begin{myquote}  #1  \end{myquote}  }


% These are usually 'schrift' not 'GroupWBlue1'
\usepackage{hyperref}   % clickable links into or outfrom this text
\hypersetup{
    colorlinks=true,
    linkcolor=GroupWBlue1,
    filecolor=magenta,      
    urlcolor=cyan,
    citecolor=myTeal
}

\numberwithin{equation}{section} % number equations as (sec.eq), not just (eq)

% Good reference on tables in LaTeX:
% https://en.wikibooks.org/wiki/LaTeX/Tables 

\usepackage{graphicx}
% parameters for the 'figure box'
\fboxsep=1mm %padding thickness
\fboxrule=2pt %border thickness

\usepackage{sectsty}

% These are usually 'schrift' not 'GroupWBlue1' or 'GroupWGrey1'
\sectionfont{\color{GroupWBlue1}}  % sets colour of sections
\subsectionfont{\color{GroupWBlue1}}
\subsubsectionfont{\color{GroupWBlue1}}

\usepackage{lastpage}

% The following two are needed to shade individual cells,
% without the `feature' of over-writing the cell boundaries.
% However, they seem to clash with empheq, so they are commented out.
% this means that \cellcolor{color}{stuff} cannot be used; see revisions just prior to 2022-03-27 for colored cells.
%\usepackage[table]{xcolor}% http://ctan.org/pkg/xcolor, https://tex.stackexchange.com/questions/50349/color-only-a-cell-of-a-table
%\usepackage{hhline}% http://ctan.org/pkg/hhline, https://tex.stackexchange.com/questions/25062/cell-shading-without-removing-lines


\usepackage[T1]{fontenc} % necessary for babel
\usepackage{babel}
\usepackage{url} % so bitex can handle URL


%\rhead{Basic CP}   % default rhead is section
%\lhead{B. P. Wise} % default lhead is subsection

% Some people say that ragged right edges make it easier to read.
\raggedright

% ---------------------------------------------
% Copyright Ben Paul Wise. All rights reserved.
% ---------------------------------------------
